\begin{abstract}
Scientific publishing systematically filters out negative results.
We argue that this long-standing asymmetry has become an urgent
problem in the era of large language models, which inherit the
positive bias of the literature they are trained on, face an
impending shortage of high-quality training data, and are
increasingly deployed as both research tools and peer reviewers. We
analyze three ways in which LLMs have changed the value of failure
data and show that the systematic absence of such data degrades
their utility as research tools, training data consumers, and peer
reviewers alike. We outline experimental protocols to validate
these claims and discuss the structural conditions under which a
failure-inclusive publishing culture could emerge.
\end{abstract}
 